%MB9T5
\loai{Ví dụ mẫu}
\begin{vidu}%KNTT
Một sợi dây đàn hồi dài 0,6 m hai đầu cố định đang có sóng dừng với một bụng sóng.
\begin{itemize}
\item[1.] Bước sóng $\lambda$ của sóng trên dây là .........
\item[2.] Giữ nguyên chiều dài sợi dây. Nếu điều chỉnh sao cho dây dao động với 3 bụng sóng thì bước sóng là .........
\item[3.] Giữ nguyên chiều dài sợi dây. Điều chỉnh dao động trên dây sao cho bước sóng đo được trên dây là 30 cm. 
\begin{itemize}
\item Số bụng sóng là ......... 
\item Số nút sóng là ......... 
\item Số bó sóng  là ......... 
\end{itemize}
\end{itemize}
\loigiai{
\begin{itemize}
\item[1.] Bước sóng: $\ell=\dfrac{\lambda}{2}\Rightarrow\lambda=2\ell=2.0,6=1,2\ m$.
\item[2.] Nếu dây dao động với 3 bụng sóng thì bước sóng là: $\ell=3\dfrac{\lambda}{2}\Rightarrow\lambda=\dfrac{2\ell}{3}=0,4\ m$.
\item[3.] Ta có: $\ell=k\dfrac{\lambda}{2}\Rightarrow k=\dfrac{2\ell}{\lambda}=8 \Rightarrow$ Số bụng $= 8$, số nút bằng 9.
\end{itemize}
}
\end{vidu}
\begin{vidu}%[3P2K8-2][Loai1]
Sóng dừng trên sợi dây đàn hồi AB (một đầu cố định, một đầu tự do), chiều dài dây là 2 m. Biết khoảng thời gian giữa hai lần sợi dây duỗi thẳng liên tiếp là 0,01 s, tốc độ truyền sóng trên dây nằm trong khoảng 75 m/s đến 85 m/s. 
\begin{itemize}
\item[1.] Chu kì sóng  là ......... ; tần số sóng là ......... 
\item[2.] Tốc độ truyền sóng là ......... 
\item[3.] Bước sóng của sóng trên dây là ......... 
\item[4.] Số bụng sóng  trên dây  là ......... ; số nút sóng trên dây  là ......... ; số bó sóng trên dây là ......... 
\end{itemize}
\loigiai{
\begin{itemize}
\item [1.] Thời gian giữa hai lần sợi dây duỗi thảng liên tiếp là $\dfrac{T}{2}$ $\Rightarrow T=0,02\ s\Rightarrow f=\dfrac{1}{T}=50\ Hz$
\item [2.] Ta có: $f=\left( 2k+1 \right)\dfrac{v}{4\ell }\Rightarrow 75\ m/s<v=\dfrac{4f\ell }{2k+1}<85\ m/s \Rightarrow 2,9 < k < 3,2 \Rightarrow k = 2\Rightarrow v = 80$ m/s.
\item [3.] Bước sóng: $\lambda=\dfrac{v}{f}=1,6\ m$
\item [4.] Ta có: $k=2$ nên số bụng = số nút = $k+1=3$
\end{itemize}
}
\end{vidu} %\dotlineEX{1}
\loai{Khoảng cách nút - bụng}
\begin{vidu}%[3P2-8.2Y][Loai1]
\nguon{MH - 2019} Trên một sợi dây đang có sóng dừng. Biết sóng truyền trên dây có bước sóng 30 cm. Khoảng cách ngắn nhất từ một nút đến một bụng là
\choice
{15 cm}
{30 cm}
{\True 7,5 cm}
{60 cm}
\loigiai{
\begin{itemize}
\item Khi có sóng dừng trên dây, khoảng cách ngắn nhất từ một nút đến một bụng là $d_{\min}=\dfrac{\lambda}{4}=7,5$ cm.
\end{itemize}
}
\end{vidu} %\dotlineEX{2}
\loai{Điều kiện sóng dừng - Tìm bước sóng, tốc độ}
\begin{vidu}%[3P2-8.2B][Loai1]
\nguon{QG - 2019} Một sợi dây dài 60 cm có hai đầu A và B cố định. Trên dây đang có sóng dừng với 2 nút sóng không kể A và B. Sóng truyền trên dây có bước sóng là
\choice
{30 cm}
{\True 40 cm}
{90 cm}
{120 cm}
\loigiai{
\begin{itemize}
\item Sóng dừng hai đầu cố định với 3 bụng: $\ell =k\dfrac{\lambda }{2}\Rightarrow \lambda =\dfrac{2\ell }{k}=\dfrac{2.60}{3}=40\ cm$.
\end{itemize}
}
\end{vidu} %\dotlineEX{2}
\loai{Hai đầu cố định - Tìm số nút, bụng}
\begin{vidu}%[3P2-8.2K][Loai1]
Trong thí nghiệm về sóng dừng, trên một sợi dây đàn hồi dài 90 cm với hai đầu cố định, tốc độ truyền sóng trên dây là 15 m/s. Biết khoảng thời gian giữa hai lần liên tiếp sợi dây duỗi thẳng là 0,02 s. Kể cả hai đầu dây, trên dây có
\choice
{7 nút và 6 bụng}
{5 nút và 4 bụng}
{\True 4 nút và 3 bụng}
{6 nút và 5 bụng}
\loigiai{
\begin{itemize}
\item Ta có: $T = 0,04\ s \Rightarrow f = 25\ Hz = k. \dfrac{v}{2\ell }\Rightarrow k = 3$ (bụng) và 4 nút trên dây.
\end{itemize}
}
\end{vidu} %\dotlineEX{2}

\loai{Một đầu cố định - Tìm số bó sóng nguyên}
\begin{vidu}%[3P2-8.2K][Loai1]
Một sợi dây đàn hồi dài 0,7 m có một đầu tự do, đầu kia nối với một nhánh âm thoa rung với tần số 80 Hz. Vận tốc truyền sóng trên dây là 32 m/s, trên dây có sóng dừng. Số bó sóng nguyên hình thành trên dây là
\choice
{6}
{\True 3}
{5}
{4}
\loigiai{
\begin{itemize}
\item Bước sóng: $\lambda = 0,4\ m\Rightarrow 0,7 = 3 \dfrac{\lambda}{2} + \dfrac{\lambda}{4}\Rightarrow$ Trên dây có 3 bó nguyên.
\end{itemize}
}
\end{vidu} %\dotlineEX{2}
\loai{Một đầu cố định - Tìm số nút, bụng}
\begin{vidu}%[3P2-8.2B][Loai1]
Một sợi dây đàn hồi dài 130 cm, có đầu A cố định, đầu B tự do dao động với tần số 100 Hz, tốc độ truyền sóng trên dây là 40 m/s. Trên dây có bao nhiêu nút và bụng sóng?
\choice
{6 nút sóng và 6 bụng sóng}
{7 nút sóng và 6 bụng sóng}
{\True 7 nút sóng và 7 bụng sóng}
{6 nút sóng và 7 bụng sóng}
\loigiai{
\begin{itemize}
\item Sóng dừng 1 đầu cố định 1 đầu tự do: $f=\left( 2n-1 \right)\dfrac{v}{4\ell }\Rightarrow $ số bụng sóng: $n = 7$ $\Rightarrow$  Số nút cũng là 7.
\end{itemize}
}
\end{vidu} %\dotlineEX{2}
\loai{Số nút bụng trên một đoạn chiều dài dây bất kì}
\begin{vidu}%[3P2-8.2K][Loai1]
Trên một sợi dây đàn hồi dài có sóng dừng với bước sóng 1,2 cm. Trên dây có hai điểm A và B cách nhau 6,1 cm, tại A là một nút sóng. Số nút sóng và bụng sóng trên đoạn dây AB là
\choice
{11 bụng, 11 nút}
{\True 10 bụng, 11 nút}
{10 bụng, 10 nút}
{11 bụng, 10 nút}
\loigiai{
$AB=6,1\ cm=10\times 0,6+0,1=10.\dfrac{\lambda}{2}+0,1\ cm \Rightarrow \begin{cases}
sb=10\\
sn=11\\
\end{cases}$
\begin{center}
\begin{tikzpicture}[scale=1,thick,>=stealth',x=1.4cm]
\draw [blue, line width = 1pt, domain=0:12, samples=500]%
plot (\x, {-0.6*cos(\x*180+90)});
\draw [blue, dashed,line width = 1pt, domain=0:12, samples=500]%
plot (\x, {0.6*cos(\x*180+90)});
\draw[red,dashed,line width = 1pt](0,0)--(12,0);
\foreach  \x  in  {0,1,2,3,4,5,6,7,8,9,10,11,12}
{\draw[dashed](\x,0)--(\x,1.4);}
\foreach  \x  in  {0,1,2,3,4,5,6,7,8,9,10,11}
{\draw[dashed,<->](\x,1.2)--++(0:1);}
\foreach  \x  in  {1,2,3,4,5,6,7,8,9,10,11,12}
{\node at (\x-0.5,-1){$\x$};}
\node[below] at (0.5,1.2){$\lambda/2$};
\draw(0,-1)--(0,2)(10.3,-1)--(10.3,2);
\draw[<->](0,1.8)--node[midway,above]{$6,1\ cm$}(10.3,1.8);
\node[above] at (0,2){$A$};
\node[above] at (10.3,2){$B$};
\end{tikzpicture}
\end{center}
}
\end{vidu} %\dotlineEX{3}
\loai{Hai đầu cố định. Cho số nút. Tìm tốc độ truyền sóng}
\begin{vidu}%[3P2-8.2K][Loai1]
\nguon{QG - 2017} Một sợi dây đàn hồi dài 90 cm có hai đầu cố định, đang có sóng dừng. Kể cả hai đầu dây cố định, trên dây có 8 nút. Biết rằng khoảng thời gian giữa 6 lần liên tiếp sợi dây duỗi thẳng là 0,25 s. Tốc độ truyền sóng trên dây là
\choice
{\True 2,6 m/s}
{1,2 m/s}
{2,9 m/s}
{2,4 m/s}
\loigiai{
\begin{itemize}
\item Khi có sóng dừng trên dây với hai đầu cố định với 8 nút ứng với 7 bó sóng $n=7$.
\item Khoảng thời gian giữa 6 lần liên tiếp sợi dây dũi thẳng ứng với $t=5\dfrac{T}{2}=0,25\Rightarrow T=0,1\ s\Rightarrow f=10\ Hz$.
\item Điều kiện để có sóng dừng trên dây $\ell=k\dfrac{v}{2f}\Rightarrow v=\dfrac{2\ell f}{k}=\dfrac{2.0,9.10}{7}=2,6$ m/s.
\end{itemize}}
\end{vidu} %\dotlineEX{3}
\loai{Một đầu cố định. Cho số bụng. Tìm tốc độ truyền sóng}
\begin{vidu}%[3P2-8.2B][Loai1]
Trên một sợi dây đàn hồi dài 1,05 m với đầu A cố định, đầu B tự do đang có sóng dừng, tần số sóng là 100 Hz. Nếu không kể đầu B người ta thấy trên dây có 3 bụng sóng. Tốc độ truyền sóng trên dây là
\choice
{40 m/s}
{35 m/s}
{50 m/s}
{\True 60 m/s}
\loigiai{
\begin{itemize}
\item Ta có: $\ell = 3,5\dfrac{\lambda}{2} \Rightarrow \lambda = \dfrac{2\ell}{3,5} = \dfrac{2.1,05}{3,5} = 0,6\ m\Rightarrow v = \lambda.f = 0,6.100 = 60\ Hz.$
\end{itemize}
}
\end{vidu} %\dotlineEX{3}
\loai{Tìm bề rộng của một bụng sóng}
\begin{vidu}%[3P2-8.2K][Loai1]
Một sóng cơ truyền trên một sợi dây rất dài thì một điểm M trên sợi dây có vận tốc dao động biến thiên theo phương trình:  $v_M=20\pi \cos(10\pi t + \varphi )$ cm/s. Giữ chặt một điểm trên dây sao cho trên dây hình thành sóng dừng và bỏ qua sự phản xạ qua lại nhiều lần. Khi đó bề rộng một bụng sóng có độ lớn là
\choice
{\True 8 cm}
{6 cm}
{16 cm}
{4 cm}
\loigiai{
\begin{itemize}
\item Ta có phương trình vận tốc:
$v = 20\pi.\cos(10\pi.t + \varphi)\Rightarrow A = 2\ cm $.
\item Bề rộng của bụng sóng là $4A=8$ cm.
\end{itemize}
}
\end{vidu} %\dotlineEX{3}
\loai{Thay đổi chiều dài dây}
\begin{vidu}%[3P2-8.2K][Loai1]
Dây căng ngang hai đầu cố định với chiều dài $\ell $, trên dây có sóng dừng. Nếu tăng chiều dài của dây lên gấp đôi (hai đầu vẫn cố định) thì dây có 10 bụng sóng, nếu tăng chiều dài thêm 30 cm (hai đầu vẫn cố định) thì trên dây có 8 nút sóng. Biết tần số, tốc độ sóng trên dây không đổi trong quá trình thay đổi chiều dài dây. Chiều dài ban đầu $\ell$ của dây là
\choice
{50 cm}
{\True 75 cm}
{150 cm}
{100 cm}
\loigiai{
\begin{itemize}
\item Ban đầu: $f=n\dfrac{v}{2\ell }$.
.\item Tăng chiều dài gấp đôi: $f=10\dfrac{v}{4\ell }$.
\item Tăng chiều dài thêm 30 cm: $f=7\dfrac{v}{2\left( \ell +30 \right)} \Rightarrow \ell  = 75$ cm.
\end{itemize}
}
\end{vidu} %\dotlineEX{6}

\begin{vidu}%[3P2-8.2K][Loai1]
Sóng dừng trên dây dài 1 m với vật cản cố định, tần số 80 Hz. Vận tốc truyền sóng là 40 m/s. Cho các điểm $M_1,\ M_2,\ M_3,\ M_4$ trên dây cách vật cản cố định là 20 cm, 30 cm, 70 cm, 75 cm. Điều nào sau đây mô tả \textbf{không} đúng trạng thái dao động của các điểm?
\choice
{$M_4$ không dao động}
{\True $M_2$ và $M_3$ dao động cùng pha}
{$M_1$ và $M_2$ dao động ngược pha}
{$M_3$ và $M_1$ dao động cùng pha}
\loigiai{
\begin{itemize}
\item $\lambda = \dfrac{v}{f} = \dfrac{40}{80} = 0,5\ m\Rightarrow$ chiều dài mỗi bó sóng là 25 cm.
\item Ta có hình vẽ\\
\begin{center}
\begin{tikzpicture}[scale=1,thick,>=stealth',x=1cm]
\tkzDefPoints{3/0/m4, 3.5/0/m3, 8.5/0/m2, 9.5/0/m1}
\draw [very thick, domain=1.5:12, samples=500]%
plot (\x, {2*cos(\x*60+90)});
\draw [very thick, domain=1.5:12, samples=500]%
plot (\x, {2*cos(\x*60-90)});
\draw(1.5,0)--(12,0);
\fill[pattern=north east lines] (12,-0.5)--(12,0.5)--(12.15,0.5)--(12.15,-0.5);
\draw(12,-0.5)--(12,0.5);
\tkzDrawPoints[size=3](m1,m2,m3,m4)
\node[above=5pt] at (m1){$M_1$};
\node[above=5pt] at (m2){$M_2$};
\node[above right=5pt] at (m3){$M_3$};
\node[above=5pt] at (m4){$M_4$};
\node[below=10pt] at (3,0){$75$};
\node[below=10pt] at (6,0){$50$};
\node[below=10pt] at (9,0){$25$};
\end{tikzpicture}
\end{center}
\item Từ hình vẽ, suy ra: $M_1,\ M_2$ ngược pha; $M_2,\ M_3$ ngược pha; $M_1,\ M_3$ cùng pha; $M_4$ đứng yên.
\end{itemize}
}
\end{vidu} %\dotlineEX{3}
